%% =====================================================================
%%  Corrigés — Unit 16 « My plans »
%% =====================================================================
\corrigesUnite{16}{My plans}

\begin{corrige}{1}
\begin{reponses}[4]
\item \rep{a trip}
\item \rep{rent a car}
\item \rep{a suitcase}
\item \rep{fly}
\item \rep{a vacation}
\item \rep{pack}
\item \rep{the beach}
\item \rep{soon}
\end{reponses}
\retiens{\en{a vacation} (américain) = \en{a holiday} (britannique) : « des vacances ».}
\end{corrige}

\begin{corrige}{2}
\begin{reponses}[2]
\item \rep{is going to fly}
\item \rep{are going to visit}
\item \rep{are going to rent}
\item \rep{isn't going to work}
\item \rep{aren't going to take}
\item \rep{Are you going to travel}
\item \rep{is Sam going to do}
\end{reponses}
\pourquoi{seul \en{be} s'accorde avec le sujet ; \en{going to} et le verbe ne changent pas. À l'oral, on
contracte : \en{We're going to rent a car}.}
\end{corrige}

\begin{corrige}{3}
\begin{reponses}[1]
\item \rep{Yes, they are.}
\item \rep{No, they aren't.}
\item \rep{No, she isn't.}
\item \rep{At Mrs.\ Lopez's apartment.}
\item \rep{In Times Square.}
\end{reponses}
\pourquoi{ses parents arrivent à 17 heures (\en{5 p.m.}) : l'après-midi, pas le matin. Le patin à glace,
c'est le 21 décembre.}
\end{corrige}

\begin{corrige}{4}
\begin{reponses}[1]
\item \rep{Are you going to travel this summer?}
\item \rep{Is Emma going to fly to London?}
\item \rep{What are you going to do tomorrow?}
\item \rep{Where are they going to stay?}
\item \rep{Is Jack going to rent a car?}
\end{reponses}
\pourquoi{comme au présent en \en{-ing}, \en{be} passe devant le sujet ; le mot interrogatif se place
en tête.}
\end{corrige}

\begin{corrige}{5}
\begin{reponses}[1]
\item \rep{I'm going to travel.}
\item \rep{She's going to fly next summer.}
\item \rep{We're not going to work.}
\item \rep{Are you going to stay?}
\item \rep{See you in two weeks!}
\end{reponses}
\retiens{\en{We're not going to…} et \en{We aren't going to…} sont tous les deux corrects.}
\end{corrige}

\begin{corrige}{6}
\begin{reponses}[3][french]
\item \rep{F}
\item \rep{V}
\item \rep{V}
\item \rep{F}
\item \rep{F}
\item \rep{V}
\end{reponses}
\pourquoi{Jack part de Chicago, à l'est (à droite sur la carte), et arrive à Los Angeles. Il dort dans
des motels et campe ; le voyage dure trois semaines (\en{three weeks}).}
\end{corrige}

\begin{corrige}{7}
\begin{reponses}[3][french]
\item \rep{plus tard}
\item \rep{maintenant}
\item \rep{maintenant}
\item \rep{plus tard}
\item \rep{plus tard}
\item \rep{maintenant}
\end{reponses}
\pourquoi{\en{going to} + verbe annonce un projet ; \en{be} + verbe en \en{-ing} seul décrit ce qui se
passe en ce moment. Texte de la piste page~\pageref{trans-16.10}.}
\end{corrige}

\begin{corrige}{8}
\begin{reponses}[4][french]
\item \rep{c}
\item \rep{a}
\item \rep{d}
\item \rep{b}
\end{reponses}
\pourquoi{personne ne va à la plage en Floride : la réponse e est en trop. Texte de la piste
page~\pageref{trans-16.11}.}
\end{corrige}

\begin{corrige}{9}
\begin{reponses}[4]
\item \rep{going}
\item \rep{visit}
\item \rep{long}
\item \rep{weeks}
\item \rep{fly}
\item \rep{stay}
\item \rep{aren't}
\item \rep{hotel}
\end{reponses}
\retiens{\en{How long…?} = « combien de temps… ? » ; \en{fly back} = « rentrer en avion ».}
\end{corrige}

\begin{corrige}{10}
\begin{reponses}[1]
\item \rep{What are you going to do this weekend?}
\item \rep{I'm going to visit my parents.}
\item \rep{We aren't going to fly.} \textit{ou} \rep{We're not going to fly.}
\item \rep{Are you going to rent a car? — Yes, I am.}
\item \rep{Next summer, they're going to go camping.}
\item \rep{She's going to arrive in two weeks.}
\item \rep{Look at those clouds! It's going to rain.}
\item \rep{One day, I'm going to see the Grand Canyon.}
\end{reponses}
\piege{phrase 2 : « rendre visite à » se dit \en{visit}, sans préposition : \en{visit my parents}.}
\end{corrige}

\begin{corrige}{11}
Réponse modèle :
\begin{anglais}
\begin{tcolorbox}[otBase,colback=otGreenTint,colframe=otGreenTint,fontupper=\fontsize{10.4}{14}\selectfont]
Next year, I'm going to travel to Canada with my sister. We're going to visit Montreal and Quebec
City. I'm going to learn to swim, and I'm going to read more books. I'm not going to watch TV every
evening. And one day, I'm going to drive on Route 66!
\end{tcolorbox}
\end{anglais}
\end{corrige}

\begin{corrige}{12}
Réponses modèles :
\begin{reponses}[1]
\item \rep{I'm going to see my friends and I'm going to cook.}
\item \rep{Yes, I am. I'm going to go to Spain.} \textit{ou} \rep{No, I'm not.}
\item \rep{I'm going to take the train.}
\item \rep{Yes, I am.} \textit{ou} \rep{No, I'm not.}
\item \rep{Sure! I'd love to!}
\end{reponses}
\end{corrige}

\begin{corrigetest}
\begin{reponses}[5][french]
\item \rep{c}
\item \rep{a}
\item \rep{b}
\item \rep{b}
\item \rep{c}
\item \rep{a}
\item \rep{c}
\item \rep{b}
\item \rep{a}
\item \rep{c}
\end{reponses}
\pourquoi{question 3 : la réponse courte reprend \en{be} (\en{isn't}), pas \en{do}. Question 5 :
\en{two weeks ago} veut dire « il y a deux semaines » (unité 17).}
\end{corrigetest}
